%% kapa: exam=1 round=25 year=2014 period=2 form=A subject=appraisal_law subject_ko=부동산관계법규 seq=25 number=25
%% source: Q-Net 감정평가사 기출문제 (cst00309) artlSeq=5101207 file=extracted/1008/5101207/2014년 제25회 감정평가사 1차시험 기출문제/2교시(A형).hwp
%% conversion: kordoc 4.24.51 -> tools/gen1.py (auto); review=proofread-sonnet; flags=none
\begin{question}{25}
국유재산법령상 일반재산에 관한 설명으로 옳은 것은?
\begin{choices}
  \item 조림을 목적으로 하는 일반재산인 토지는 20년간 대부할 수 있다.
  \item 증권을 제외한 일반재산을 처분할 때에 그 대장가격이 3천만원 미만인 경우에는 두 개의 감정평가법인의 평가액을 산술평균한 금액으로 예정가격을 결정한다.
  \item 일반재산인 토지와 사유재산인 토지를 교환할 때 쌍방의 가격이 같지 아니하면 그 차액을 금전, 증권 또는 현물로 대납(代納)할 수 있다.
  \item 일반재산을 매각하는 경우 매수자에게 그 재산의 용도와 그 용도에 사용하여야 할 기간을 정하여 매각하는 것은 허용되지 않는다.
  \item 중앙관서의 장이 소관 특별회계에 속하는 일반재산 중 일단의 토지면적이 4천제곱미터인 재산을 매각하려는 경우에는 총괄청과 협의하여야 한다.
\end{choices}
\end{question}
\answer{5}
